Les piles ou accumulateurs exploitent l'énergie électrique qu'elles fournissent
pour faire fonctionner plusieurs appareils électriques. Cette partie vise à
étudier un exemple de ces piles~: \textbf{la pile zinc/nickel}.

Pour réaliser la pile zinc/nickel, au cours d'une séance de travaux pratiques,
un groupe d'élèves a utilisé les outils et solutions suivants~:
\begin{itemize}
    \item Un bécher contenant un volume $V_1=\qty{20}{\milli\litre}$ d'une
          solution aqueuse de nitrate de nickel~:
          ${\ce{Ni^{2+}_{(aq)} + 2NO^{-}_{3(aq)}}}$ de concentration molaire~:
          $C_1=\qty{1,0e-1}{\mole\per\litre}$
    \item Un bécher contenant un volume $V_2=\qty{20}{\milli\litre}$ d'une
          solution aqueuse de sulfate de zinc~:
          ${\ce{Zn^{2+}_{(aq)} + SO^{2-}_{4(aq)}}}$ de concentration molaire~:
          $C_2=\qty{5,0e-2}{\mole\per\litre}$
    \item Un fil de zinc et un autre de nickel~;
    \item Un pont salin.
\end{itemize}

\begin{donnees}
    \begin{itemize}
        \item ${M(\ce{Zn})=\qty{65,4}{\gram\per\mole}}$~; \
              ${1\mathcal{F}=\qty{96500}{\coulomb\per\mole}}$
    \end{itemize}
\end{donnees}
Un élève a réalisé le circuit électrique suivant en utilisant la pile
zinc/nickel, un ampèremètre et un conducteur ohmique. Il a observé, après la
fermeture du circuit, le passage d'un courant électrique dans l'ampèremètre,
dirigé de l'électrode de nickel vers l'électrode de zinc à l'extérieur de la
pile, et d'intensité constante $I$.

\begin{enumerate}
    \item Donner la représentation conventionnelle de la pile.~\textbf{(1pt)}
    \item Écrire l'équation chimique modélisant la transformation ayant lieu
          pendant le fonctionnement de la pile.~\textbf{(1pt)}
    \item Après une durée de fonctionnement $\Delta t=\qty{2}{\hour}$, la pile
          est devenue usée.
          \begin{enumerate}
              \item Établir le tableau d'avancement décrivant l'évolution du
                    système chimique.~\textbf{(1pt)}
              \item Déterminer le réactif limitant, sachant que la masse de la
                    partie immergée du fil de zinc est~:
                    $m=\qty{1,0}{\gram}$.~\textbf{(1pt)}
              \item Calculer la valeur de l'intensité $I$.~\textbf{(1,5pt)}
          \end{enumerate}
\end{enumerate}

